\documentclass[10pt]{beamer}

\usepackage{verbatim}
%\usepackage[latin1]{inputenc}
\usepackage{amssymb,amsmath,amsfonts,url}
\usepackage{graphicx,color}
\usepackage{pgf,pgfarrows,pgfnodes,pgfautomata,pgfheaps,pgfshade}
\usepackage[utf8]{inputenc}
\usepackage[french]{babel}

\usepackage[all]{xy}
\usepackage{amstext}
\usepackage{xspace}
\usepackage{mdwtab}


\mode<article> {
  \usepackage{fullpage}
  \usepackage{hyperref}
  \setjobnamebeamerversion{advancedcourse.beamer}
}




\usepackage{listings}
\usepackage{xcolor}

\lstset{
	language=Python,
	basicstyle=\ttfamily\small,
	keywordstyle=\color{blue},
	commentstyle=\color{gray},
	stringstyle=\color{green!50!black},
	showstringspaces=false,
	frame=single,
	numbers=none
}
\lstset{
	language=Python,
	basicstyle=\ttfamily\small,
	keywordstyle=\color{blue},
	commentstyle=\color{gray},
	stringstyle=\color{green!50!black},
	showstringspaces=false,
	frame=single,
	numbers=none,
	literate=
	{é}{{\'e}}1
	{è}{{\`e}}1
	{ê}{{\^e}}1
	{ë}{{\"e}}1
	{à}{{\`a}}1
	{â}{{\^a}}1
	{ä}{{\"a}}1
	{ù}{{\`u}}1
	{û}{{\^u}}1
	{ü}{{\"u}}1
	{ô}{{\^o}}1
	{ö}{{\"o}}1
	{î}{{\^i}}1
	{ï}{{\"i}}1
	{ç}{{\c c}}1
}


\newcommand{\alertonce}[1]{\addtocounter{beamerpauses}{-1}\alert<+>{#1}}
\newcommand{\alerttwice}[1]{\addtocounter{beamerpauses}{-1}\alert<+-+(1)>{#1}}

%\newcommand{\subsubsubsection}[1]{\noindent\textbf{#1}}

\newcommand{\cfbox}[1]{\begin{center}\fbox{#1}\end{center}}
\newcommand{\cbox}[1]{\begin{center} #1 \end{center}}

%%%%%%%%%%%%%%%%%%%%%%%%%%%%%%%%%%%%%%%%%%%%%%%%%%%%%%%%%%%%%%%%%%%%%%%%%%%%%%%%%%%%%%%%%

\graphicspath{{Graphics/}}
%\input def_0.tex

%\input{StyleSlides.tex}
\mode<presentation> {
    \usepackage{beamerthemeBoadilla}

  \beamertemplatetransparentcovereddynamic
  \beamertemplateballitem
}


\title[IN100 Cours 4]{\textbf{Fondements informatiques I}\\
Cours 4: Fonctions}

\author[]{Sorina Ionica \texttt{sorina.ionica@uvsq.fr} \\~\\Bilal FAYE  \texttt{bilal.faye@uvsq.fr} }


\institute{}

\date{}


%\setbeamercolor{normal text}{fg=black}

\begin{document}


\frame{\titlepage}


\begin{frame}
	\frametitle{Les fonctions}
	
	\begin{block}{Pourquoi utiliser une fonction ?}
		Une fonction permet de \textbf{regrouper plusieurs instructions}
		pour réaliser une tâche précise.
		
		\vspace{0.1cm}
		
		Elle peut ensuite être \textbf{réutilisée} plusieurs fois dans un programme.
	\end{block}
	
	\vspace{0.2cm}
	
	\begin{exampleblock}{Idée générale}
		\centering
		\textbf{Données}
		$\;\longrightarrow\;$
		\textbf{Fonction}
		$\;\longrightarrow\;$
		\textbf{Résultat}
	\end{exampleblock}
	
	\vspace{0.2cm}
	
	\begin{alertblock}{Exemple}
		Une fonction pourrait prendre deux nombres en entrée
		et calculer leur somme.
		
		\[
		3,\;5
		\quad\longrightarrow\quad
		\text{calcul de la somme}
		\quad\longrightarrow\quad
		8
		\]
	\end{alertblock}
	
\end{frame}
\begin{frame}
	\frametitle{Comment fonctionne une fonction ?}
	
	\begin{block}{Une fonction peut recevoir des données et produire un résultat}
		Une fonction est généralement composée de trois éléments :
	\end{block}
	
	\vspace{0.15cm}
	
	\begin{columns}[T]
		\begin{column}{0.3\textwidth}
			\begin{exampleblock}{Entrées}
				\centering
				Les données\\
				nécessaires
			\end{exampleblock}
		\end{column}
		
		\begin{column}{0.3\textwidth}
			\begin{exampleblock}{Traitement}
				\centering
				Les instructions\\
				à exécuter
			\end{exampleblock}
		\end{column}
		
		\begin{column}{0.3\textwidth}
			\begin{exampleblock}{Sortie}
				\centering
				Le résultat\\
				obtenu
			\end{exampleblock}
		\end{column}
	\end{columns}
	
	\vspace{0.2cm}
	
	\begin{alertblock}{Exemple}
		\centering
		Deux nombres
		$\;\longrightarrow\;$
		\textbf{addition}
		$\;\longrightarrow\;$
		leur somme
	\end{alertblock}
	
\end{frame}
\begin{frame}
	\frametitle{Les fonctions en mathématiques}
	
	\begin{block}{Une fonction}
		En mathématiques, une fonction associe une ou plusieurs
		\textbf{entrées} à un \textbf{résultat}.
	\end{block}
	
	\vspace{0.15cm}
	
	\begin{exampleblock}{Notation}
		\centering
		$f : x \rightarrow f(x)$
		
		\vspace{0.1cm}
		
		\textbf{$f$} : nom de la fonction
		\qquad
		\textbf{$x$} : argument
		\qquad
		\textbf{$f(x)$} : résultat
	\end{exampleblock}
	
	\vspace{0.15cm}
	
	\begin{columns}[T]
		\begin{column}{0.48\textwidth}
			\begin{exampleblock}{Exemples}
				\[
				f(x) = ax + b
				\]
				\[
				f(x,y) = x+y
				\]
			\end{exampleblock}
		\end{column}
		
		\begin{column}{0.48\textwidth}
			\begin{exampleblock}{Autres fonctions}
				\[
				f(x) = x^2
				\]
				\[
				f(x) = 2^x
				\]
				\[
				f(x) = |x|
				\]
			\end{exampleblock}
		\end{column}
	\end{columns}
	
\end{frame}
\begin{frame}[containsverbatim]
	\frametitle{Définir une fonction en Python}
	
	\begin{alertblock}{Syntaxe}
\begin{lstlisting}[language=Python]
def nom_fonction(arg1, arg2, ...):
	instruction 1
	instruction 2
	...
	insstruction n
\end{lstlisting}
	\end{alertblock}
	
	\vspace{0.1cm}
	
	\begin{itemize}
		\item \texttt{def} indique que l'on définit une fonction.
		
		\item On indique ensuite le \textbf{nom} de la fonction et ses
		\textbf{arguments} entre parenthèses.
		
		\item Le \textbf{corps de la fonction} est un bloc d'instructions
		\textbf{indenté}.
	\end{itemize}
	
	\vspace{0.1cm}
	
	\begin{exampleblock}{Lien avec les mathématiques}
		\centering
		$f(x)$
		$\quad\longleftrightarrow\quad$
		\texttt{nom\_fonction(argument)}
	\end{exampleblock}
	
\end{frame}
\begin{frame}[containsverbatim]
	\frametitle{Définir une fonction en Python}
	
	\begin{exampleblock}{Exemple}
\begin{lstlisting}[language=Python]
def multiplie(x):
	"""Multiplie x par 2 et affiche le résultat."""
	print(x * 2)
multiplie(7)
multiplie("aaa")
\end{lstlisting}
	\end{exampleblock}
	
	\vspace{0.1cm}
	
	\begin{itemize}
		\item \texttt{multiplie} est le \textbf{nom} de la fonction.
		\item \texttt{x} est son \textbf{argument}.
		\item \texttt{multiplie(7)} appelle la fonction avec \texttt{x = 7}.
	\end{itemize}
	
	\vspace{0.05cm}
	
	\begin{alertblock}{Docstring}
		La chaîne entre triples guillemets \texttt{""" ... """} décrit
		le rôle de la fonction. Elle peut notamment être consultée avec
		\texttt{help()}.
	\end{alertblock}
	
\end{frame}
\begin{frame}[containsverbatim]
	\frametitle{L'instruction \texttt{return}}
	
	\begin{block}{Rôle de \texttt{return}}
		\texttt{return} permet à une fonction de \textbf{renvoyer une valeur}
		au programme qui l'a appelée.
		
		\vspace{0.05cm}
		
		Une fonction peut ne pas avoir de \texttt{return} :
		dans ce cas, elle renvoie \texttt{None}.
	\end{block}
	
	\vspace{0.15cm}
	
	\begin{columns}[T]
		\begin{column}{0.48\textwidth}
			\begin{exampleblock}{Avec \texttt{return}}
\begin{lstlisting}[language=Python]
def multiplie(x):
	return x * 2
print(type(multiplie(2)))
\end{lstlisting}
				
				\centering
				La fonction renvoie \texttt{4}.
			\end{exampleblock}
		\end{column}
		
		\begin{column}{0.48\textwidth}
			\begin{exampleblock}{Sans \texttt{return}}
\begin{lstlisting}[language=Python]
def multiplie(x):
	print(x * 2)
print(type(multiplie(2)))
\end{lstlisting}
				
				\centering
				La fonction renvoie \texttt{None}.
			\end{exampleblock}
		\end{column}
	\end{columns}
	
	\vspace{0.1cm}
	
	\begin{alertblock}{À retenir}
		\centering
		\texttt{print()} affiche une valeur \qquad
		\textbf{vs.} \qquad
		\texttt{return} renvoie une valeur
	\end{alertblock}
	
\end{frame}
\begin{frame}[containsverbatim]
	\frametitle{Signature d'une fonction}
	
	\begin{block}{Définition}
		La \textbf{signature d'une fonction} décrit les éléments
		qui permettent de caractériser la fonction :
		\textbf{son nom, ses paramètres et leurs types},
		ainsi que le \textbf{type de la valeur renvoyée}.
	\end{block}
	
	\vspace{0.2cm}
	
	\begin{alertblock}{En Python}
		Les types des paramètres et le type de retour
		\textbf{ne sont pas obligatoires} dans la définition.
		
		\medskip
		
		Python détermine les types lors de l'exécution.
	\end{alertblock}
	
	\vspace{0.15cm}
	
	\begin{exampleblock}{Que se passe-t-il lors de l'appel ?}
		\begin{lstlisting}[language=Python]
def somme(x, y):
	return x + y
print(somme("hello", 1))
\end{lstlisting}
		

		%Python essaie d'effectuer \texttt{"hello" + 1}
		%et provoque une \textbf{erreur de type}.
	\end{exampleblock}

	
\end{frame}
\begin{frame}[containsverbatim]
	\frametitle{Retourner plusieurs valeurs}
	\begin{block}{Une fonction peut renvoyer plusieurs valeurs}
		En Python, une instruction \texttt{return} peut renvoyer
		plusieurs valeurs séparées par des virgules.
		
		\vspace{0.05cm}
		
		Ces valeurs sont regroupées dans un \textbf{tuple}.
	\end{block}
	
	\vspace{0.15cm}
	
	\begin{exampleblock}{Exemple}
		\begin{lstlisting}[language=Python]
def ordonne(a, b):
	if a < b:
		return a, b
	else:
		return b, a
print(ordonne(10, 5))
		\end{lstlisting}
	\end{exampleblock}
	

	\vspace{0.1cm}
	
	\begin{alertblock}{Résultat}
		\centering
		\texttt{(5, 10)}
		
		\vspace{0.05cm}
		
		La fonction renvoie donc un \textbf{tuple contenant deux valeurs}.
	\end{alertblock}
	
	
\end{frame}

\begin{frame}[containsverbatim]
	\frametitle{Sortir de la fonction}

		\small Dès qu'une fonction atteint une instruction \texttt{return},
		elle \textbf{s'arrête immédiatement}. Les instructions placées après ce \texttt{return}
		ne sont donc pas exécutées.

	

	
	\begin{exampleblock}{Exemple}
\begin{lstlisting}[language=Python]
def un():
	return 1
	print("Je ne serai jamais atteint")
print(un())
\end{lstlisting}
	\end{exampleblock}
	


	
	\begin{alertblock}{Dans une boucle}
		\begin{lstlisting}[language=Python]
def cherche(x, l):
	for elem in l:
		if x == elem:
			return 1
	return 0
print(cherche(10, [1, 10, 0]))  # 1
		\end{lstlisting}
		
		
		\small Dès que \texttt{10} est trouvé, \texttt{return 1} termine
		la fonction : les éléments suivants ne sont pas parcourus.
	\end{alertblock}
	
	
\end{frame}
\begin{frame}[containsverbatim]
	\frametitle{Les fonctions existantes en Python}
	

	\begin{block}{Pas besoin de tout programmer soi-même}
		Python fournit de nombreuses \textbf{fonctions déjà définies}.
		On peut les appeler directement dans notre programme.
	\end{block}
	
	\vspace{0.15cm}
	
	\begin{exampleblock}{Quelques fonctions que vous connaissez déjà}
\begin{lstlisting}[language=Python]
print("Bonjour")       # affiche une valeur
type(3.14)             # donne le type de la valeur
len("Bonjour")         # donne la longueur
\end{lstlisting}
	\end{exampleblock}
	

	\vspace{0.1cm}
	
	\begin{alertblock}{À retenir}
		Une fonction existante s'utilise comme une fonction que
		l'on a définie nous-mêmes : \textbf{nom(arguments)}.
	\end{alertblock}
	
	
\end{frame}

\begin{frame}[containsverbatim]
	\frametitle{Les fonctions d'un module}
	\begin{block}{Certaines fonctions sont regroupées dans des modules}
		Un \textbf{module} est un ensemble de fonctions et d'outils
		déjà programmés que l'on peut utiliser.
		
		\vspace{0.05cm}
		
		Pour utiliser un module, on commence par l'\textbf{importer}.
	\end{block}
	
	\vspace{0.15cm}
	
	\begin{exampleblock}{Exemple avec le module \texttt{math}}
		\begin{lstlisting}[language=Python]	
import math
print(math.sqrt(25))
		\end{lstlisting}
	\end{exampleblock}
	

	\vspace{0.1cm}
	
	\begin{alertblock}{Comment lire \texttt{math.sqrt(25)} ?}
		\texttt{math} : le module \qquad
		\texttt{sqrt} : la fonction
		
		\[
		\texttt{math.sqrt(25)} = 5.0
		\]
	\end{alertblock}

	
\end{frame}
\begin{frame}[containsverbatim]
	
	
	\begin{exampleblock}{Que fait ce programme ?}
		\begin{lstlisting}[language=Python]
import math
def calcule(x):
	if x > 0:
		return math.sqrt(x)
	else:
		return 0
i = 1
while i <= 5:
	for x in [i, i + 1]:
		print(calcule(x))
	i = i + 1
		\end{lstlisting}
	\end{exampleblock}
\end{frame}






































\begin{comment}

\frame[containsverbatim]{
\frametitle{Les fonctions}


\begin{itemize}
\item Il est très fréquent d'utiliser la même
fonctionnalité plusieurs fois.
\item La librairie standard de Python ou des librairies externes
vont vous fournir des fonctions déj\`a codées.
\end{itemize}

\begin{exampleblock}{Exemples}
\begin{verbatim}
import math

x=2
math.sqrt(x)
math.cos(x)
print(x)
chaine="hello"
x=len(chaine)
\end{verbatim}
\end{exampleblock}

}


\frame[containsverbatim]{
\frametitle{Les fonctions}


En mathématique une fonction s'écrit $f: x \rightarrow  f(x)$. \\
\vspace{0.2 cm}

 Il y a trois parties dans la fonction: $f$ son \emph{nom}, $x$ son \emph{argument}
et \emph{f(x)} son image (valeur/résultat).\\
\vspace{0.2 cm}

Quelques exemples que vous connaissez bien: 

\begin{columns}
\column{5 cm}
\begin{eqnarray*}
 f(x) &=& ax + b \\
 f(x,y) &= &x + y
\end{eqnarray*}

\column{5 cm}
\begin{eqnarray*}
f(x) &= & x^2\\ 
f(x)&=&2^x\\
f(x) & =& |x|\\
\end{eqnarray*}
\end{columns}
}




\frame[containsverbatim]{
\frametitle{Définir une fonction en Python}


\begin{alertblock}{Syntaxe}
\begin{verbatim}
def nom_fonction(arg 1, arg2, ..., argn)
       instruction1 
       instruction 2
       ...
       instruction n 

\end{verbatim}
\end{alertblock}


\begin{itemize}

\item  Une fonction est introduite grâce au mot clé \texttt{def}
\item
On indique ensuite son nom et on donne sa liste d'arguments entre parenthèses
\item
Le corps de la fonction est un  \emph{bloc d'instructions}, donc indenté, qui vient
après sa déclaration. 
 \end{itemize}



}

\frame[containsverbatim]{
\frametitle{Définir une fonction en Python}



\begin{exampleblock}{Un premier exemple}
\begin{verbatim}
def multiplie(x):
   """ Cette fonction multiplie son argument par 2 et l'affiche"""
   print(x*2)

multiplie(7)
multiplie("aaa")
\end{verbatim}
\end{exampleblock}

\vspace{0.5 cm}

\begin{itemize}

\item Le bloc de la fonction peut être précédé d'un commentaire, appelé
\textbf{docstring} qui sert \`a générer la documentation automatique de
\texttt{help()}. Un docstring est entouré de trois guillemets:
\texttt{"""\ docstring\ """}.
 \end{itemize}



}

\frame[containsverbatim]{
\frametitle{\'Ecrire une fonction}

\begin{itemize}
 \item Les fonctions retournent une \textbf{valeur},  donnée après le mot
clé \texttt{return}. 
\item Même les fonctions qui n'ont pas d'instruction
\texttt{return}, ou qui terminent en finissant d'exécuter le bloc de code sans
atteindre un \texttt{return}, retournent la valeur \texttt{None}.
\end{itemize}


\begin{alertblock}{Syntaxe}
\begin{verbatim}
def nom_fonction(arg 1, arg2, ..., argn):
       instruction1 
       instruction 2
       ...
       return objet

\end{verbatim}
\end{alertblock}

\begin{columns}
\column{5 cm}
\begin{exampleblock}{Avec \texttt{return}}
\begin{verbatim}
def multiplie(x):
        return x*2
        
print(type(multiplie(2)))            
\end{verbatim}
\end{exampleblock}

\column{5 cm}
\begin{exampleblock}{Sans \texttt{return}}
\begin{verbatim}
def multiplie(x):
       print(x*2)
       
print(type(multiplie(2)))      
\end{verbatim}
\end{exampleblock}
\end{columns}
}



\frame[containsverbatim]{
\frametitle{Signature d'une fonction}

\begin{itemize}
\item En programmation, la \emph{signature d'une fonction} est donnée par son nom, les paramètres, leur type et le type renvoyé. 
\item En Python, pas besoin de donner le type renvoyé lors de la déclaration, ni le type des paramètres. 
\end{itemize}

\vspace{0.5 cm}

Que se passe-t-il lors de l'appel de la fonction? 

\vspace{0.5 cm}

\begin{exampleblock}{}
\begin{verbatim}
def somme(x,y):
    return x+y
    
print(somme("hello", 1))    
\end{verbatim}

\end{exampleblock}


}




\frame[containsverbatim]{
\frametitle{Sortir de la fonction}

\begin{itemize}
 \item  Attention, une fonction termine dès qu'elle atteint le mot clé \texttt{return}. 
 \item  C'est utile pour gérer le flot d'exécution d'une fonction. 
 \end{itemize}

\begin{exampleblock}{Exemple}
\small{
\begin{verbatim}
def un():
    return 1
    print("je ne serai jamais atteint")      
un()

def cherche(x,l):
    for elem in l:
        1 / elem #erreur quand elem vaut 0
        if(x == elem):
            return 1
    return 0

cherche(10, [1,10,0]) #la valeur 0  pas atteinte car on trouve 10 avant
\end{verbatim}
}
\end{exampleblock}

}



\frame[containsverbatim]{
\frametitle{Type de return}

\begin{itemize}
\item En Python, on peut utiliser des types de retour complexes (contrairement \`a C par exemple).
\item On peut par exemple renvoyer des tuples de valeurs, ce qui est souvent pratique.
\end{itemize}

\vspace{0.5 cm}

\begin{exampleblock}{Exemple}
\begin{verbatim}
def ordonne(a, b):
    if a < b:
        return a, b
    else:
        return b, a
    
print(ordonne(10, 5))
\end{verbatim}
\end{exampleblock}

}
\end{comment}

\end{document}

